\documentclass[../Thesis]{subfiles}
\begin{document}

\chapter{Feasibility} \label{chapter4}
\chaptermark{Feasibility}
\section{Introduction}
This chapter evaluates the feasibility of the WaveID system before full implementation. It presents preliminary dataset analysis, prototype development, tool validation, and early experimental findings. The goal is to confirm that the project's core objectives are achievable within the technical, temporal, and ethical constraints identified.

\section{Preliminary Dataset Studies}
The first stage of feasibility assessment involved examining the available datasets to ensure suitability for training and evaluation.

\subsection{Dataset Quality}
Both the Free Music Archive (FMA) and GTZAN datasets were assessed for:
\begin{itemize}
    \item \textbf{Audio consistency:} Files contain minimal clipping, artefacts, or encoding noise.
    \item \textbf{Genre diversity:} A wide range of musical styles, improving generalisation.
    \item \textbf{Licensing:} Both datasets are legally distributable and appropriate for academic research.
\end{itemize}

Minor issues such as duplicated tracks and inconsistent metadata were identified but documented for later mitigation. These inconsistencies do not threaten the feasibility of the project.

\subsection{Segmenting for Short-Form Media}
To simulate short-form video excerpts, tracks were segmented into windows of 1-3 seconds. Early inspection confirmed that even short segments retained sufficient melodic or harmonic structure to support identification models.

\section{Prototype System Architecture}
A minimal prototype will be developed to validate core system components, including preprocessing, embedding extraction, and similarity matching.

\subsection{Preprocessing Validation}
The preprocessing pipeline will successfully:
\begin{itemize}
    \item Resample tracks to a uniform rate,
    \item Normalise amplitude levels,
    \item Generate mel spectrograms using \texttt{Librosa}.
\end{itemize}

This will validate that preprocessing can be carried out efficiently and at scale.

\subsection{Embedding Extraction Prototype}
Two embedding pathways that will be tested:
\begin{itemize}
    \item \textbf{Pre-trained models} such as VGGish and OpenL3,
    \item \textbf{A small contrastive-learning prototype} using pitch-shift and time-stretch augmentation.
\end{itemize}

The goal is for it to indicate that embeddings cluster semantically similar audio, confirming that the models capture meaningful acoustic relationships.

\section{Tooling and Development Environment}
A feasibility check was conducted on the software stack required. All required tools were successfully installed and validated, including Python 3.10, Librosa for audio feature extraction, PyTorch and TensorFlow for model development, and FAISS for efficient similarity search. These tools worked without dependency conflicts, confirming that the development environment is stable and suitable for the project. All tools ran without compatibility issues.

\subsection{Hardware Feasibility}
The project requires modest GPU computation. Tests showed:
\begin{itemize}
    \item Pre-trained embeddings run efficiently on CPU,
    \item Contrastive training benefits significantly from GPU support,
    \item FAISS GPU mode is optional but beneficial.
\end{itemize}

These findings confirm that the system is feasible within the available computational resources.

\section{Preliminary Experiments}
Early experiments will be conducted to validate that the intended approach can reliably identify transformed audio.

\subsection{Transformation Testing}
Early testing will examine model performance on pitch-shifted, time-stretched, cropped, and noise-augmented samples. The working hypothesis is that pre-trained embeddings will perform well on clean audio but degrade under strong transformations, while contrastive-learning embeddings will show greater resilience to these changes. These initial results will inform the final model design and guide augmentation strategy.


The current hypothesis is that:
\begin{itemize}
    \item \textbf{Pre-trained embeddings} will perform well on clean audio but will most likely degrade under stronger transformations,
    \item \textbf{Contrastive-learning embeddings} should maintain higher similarity scores under pitch and tempo modifications.
\end{itemize}

\section{Feasibility Conclusion}
Overall, the feasibility analysis supports the viability of WaveID. Preliminary experiments indicate that embedding-based identification is achievable, contrastive learning improves robustness to transformations, and all system components, preprocessing, embedding extraction, similarity search, and evaluation operate as intended in prototype form. These findings confirm that the project is ready to progress to full implementation.

\end{document}
